\begin{textsample}{NEG Dim 4 – 2023\_c5 – Score -1.00 – 2023\_c5/enade\_2023\_odonto\_f04d613ff6.txt}  \label{ex:f4_neg_094}
CAPÍTULO 4 DISTRIBUIÇÃO DE CONCEITOS A finalidade deste capítulo é apresentar um panorama nacional sobre a distribuição do Conceito Enade 2023.
Para traçar esse panorama, recorre -se a tabelas e gráficos, que são acompanhados de comentários e análises que contribuem para a contextualização dos dados neles expostos.
Para visualizar o panorama do Enade 2023, em nível nacional, na primeira seção deste capítulo, os dados sobre o número de cursos participantes são organizados por grande região.
Na segunda, o mapeamento do conceito do exame é feito por categoria administrativa das IES e modalidade de oferta dos cursos, de maneira cruzada às grandes regiões.
Na última seção, a lógica de apresentação dos dados é feita por organização acadêmica e, também, por grande região de funcionamento dos cursos, no país como um todo. 4.1 PANORAMA NACIONAL DA DISTRIBUIÇÃO DOS CONCEITOS A lógica adotada na apresentação dos dados, nesta seção e nas duas seguintes, pressupõe o registro de duas informações de natureza reiterativa.
O primeiro diz respeito aos valores possíveis do Enade, que variam em níveis de 1 a 5, sendo que, à medida que esses valores aumentam, o curso é melhor avaliado, e ao fato de que aqueles cursos que tiveram somente um ou nenhum concluinte participante não obtêm o Conceito Enade, ficando sem conceito ( SC ), como ressaltado na apresentação deste relatório.
O segundo registro é que, caso haja diferença entre os cursos tabulados neste capítulo e no capítulo 1, ela corresponderá aos cursos que não foram avaliados, ou seja, àqueles, em princípio, sem estudantes concluintes que tivessem sido inscritos para a prova.
Na Tabela 4.1, são apresentadas a quantidade e a distribuição de cursos de Odontologia participantes do Enade 2023, por faixa de conceito e grande região.
A partir dos dados da Tabela 4.1, nota -se que, dos 452 cursos participantes, 184 ( 40,7 \% ) classificaram -se com conceito 3, o valor modal.
Esse foi também o conceito modal nas regiões : Nordeste ( 35,2 \% ), Sudeste ( 44,2 \% ) e Sul ( 50,6 \% ).
O conceito 2 foi o segundo mais frequente em nível nacional ( 30,5 \%, correspondendo a 138 cursos ), tendo sido o conceito modal ou segundo conceito mais frequente na maioria das regiões.
O conceito 4 foi o terceiro ( 18,1 \%, correspondendo a 82 cursos ).
Houve, ainda, 25 ( 5,5 \% ) cursos que receberam o conceito 5 e 19 ( 4,2 \% ) cursos que receberam o conceito 1.
Dos 452 cursos de Odontologia, quatro ( 0,9 \% ) ficaram sem conceito ( SC ).
A região Norte participou com 43 cursos, o que corresponde a 9,5 \% do total nacional.
Desses, 21 cursos receberam o conceito 2, o que equivale a 48,8 \% do total regional, o conceito modal na região.
A 15 cursos ( 34,9 \% ) atribuiu -se o conceito 3 ; a cinco cursos ( 11,6 \% ), o conceito 1 ; e a um curso ( 2,3 \% ) o conceito 5.
Ainda nessa região, um curso ( 2,3 \% ) obteve conceito 4 e nenhum curso ( 0,0 \% ) ficou sem conceito ( SC ).
A região Nordeste participou com 122 cursos, correspondentes a 27,0 \% do total nacional.
Nessa região, 43 cursos ( 35,2 \% em termos regionais ) obtiveram os conceitos 3 e 2, conceitos modais para a região.
O conceito 4 foi atribuído a 22 cursos ( 18,0 \% ).
Ainda nessa região, sete cursos ( 5,7 \% ) foram avaliados com o conceito 5 e seis cursos ( 4,9 \% ) foram avaliados com o conceito 1.
Um curso ( 0,8 \% ) ficou sem conceito ( SC ).
Para os 165 cursos participantes da região Sudeste ( 36,5 \% do total nacional ), o conceito 3 foi o conceito modal, como já anteriormente assinalado, tendo sido obtido por 73 cursos ( 44,2 \% ).
O conceito 2 foi atribuído a 45 cursos ( 27,3 \% ) e o conceito 4, a 34 cursos ( 20,6 \% ).
Receberam os conceitos 5 e 1, 103 Relatório Síntese de Área – ENADE/2023 respectivamente, dez cursos ( 6,1 \% ) e dois cursos ( 1,2 \% ).
Nessa região, um curso ( 0,6 \% ) ficou sem conceito ( SC ).
Os 77 cursos da região Sul corresponderam a 17,0 \% do total nacional.
Houve predominância do conceito 3 ( conceito modal ), atribuído a 50,6 \% dos cursos da região, ou seja, atribuído a 39 dos 77 cursos participantes na região Sul.
O conceito 4 foi atribuído a 19 cursos ( 24,7 \% ) ; o conceito 2, 12 cursos ( 15,6 \% ) ; e o conceito 5, a cinco cursos ( 6,5 \% ).
Na região Sul, nenhum curso ( 0,0 \% ) obteve conceito 1 e dois cursos ( 2,6 \% ) ficaram sem conceito ( SC ).
Já dos 45 cursos participantes na região Centro-Oeste ( 10,0 \% do total nacional ), 17 ( 37,8 \% em termos regionais ) receberam conceito 2, o conceito modal. 14 cursos ( 31,1 \% ) obtiveram o conceito 3, seis cursos ( 13,3 \% ) obtiveram o conceito 4 e seis cursos ( 13,3 \% ), o conceito 1.
Nessa região, dois cursos ( 4,4 \% ) obtiveram o conceito 5 e nenhum curso ( 0,0 \% ) ficou sem conceito ( SC ).
Em um gráfico de distribuição cumulativa, em cada valor de abscissa ( eixo x ), o valor apresentado no eixo das ordenadas ( eixo y ) é o valor acumulado da variável ( Conceito Enade ).
Em linhas gerais, em um gráfico desta mesma tipologia, poligonais mais à esquerda representam, no caso dos dados apresentados nesta seção, as grandes regiões com a pior distribuição, e poligonais mais à direita, as grandes regiões com os melhores conceitos.
A partir destas diretrizes, no Gráfico 4.1, é apresentada a distribuição cumulativa do Conceito Enade, segundo a grande região.
No Gráfico 4.1, observando a linha vertical pontilhada, que ajuda a visualizar o percentual de conceitos acumulados até o valor 3, ou seja, conceitos 1, 2 e 3, pode -se notar, por exemplo, que a região Norte ( linha azul ) apresenta 95,3 \% dos cursos com conceito 3 ou abaixo deste valor.
A região Sul ( linha roxa ), por outro lado, apresenta 68,0 \% dos cursos na mesma situação : com conceito 3 ou abaixo.
A região Norte apresenta os cursos com conceitos mais baixos, ao passo que a região Sul apresenta a distribuição com valores maiores.
Os cursos em IES situadas na região Nordeste ( linha verde ) apresentam uma situação intermediária. 4.2 CONCEITOS POR CATEGORIA ADMINISTRATIVA, MODALIDADE DE OFERTA DOS CURSOS E GRANDE REGIÃO Na Tabela 4.2, é apresentada a distribuição dos cursos de Odontologia participantes do Enade 2023, por categoria administrativa e modalidade de oferta, de acordo com os conceitos por eles alcançados, segundo as grandes regiões brasileiras.
De acordo com as informações da Tabela 4.2, dos 452 cursos participantes, 61 ( 13,5 \% ) eram ministrados em instituições públicas e 391 ( 86,5 \% ), em instituições privadas.
Quanto à modalidade de oferta, todos os 452 ( 100,0 \% ) cursos eram presenciais.
De acordo com as informações da mesma tabela, em termos nacionais, entre instituições públicas e privadas, observa -se que, dos 25 cursos avaliados com conceito 5, 18 foram oferecidos em IES públicas, e 7, em IES privadas.
Dos 391 cursos participantes de IES privadas, o conceito 3 foi o valor modal, atribuído a 181 ( 46,3 \% ) cursos.
Em relação aos demais cursos de instituições privadas, 130 cursos ( 33,2 \% ) foram avaliados com o conceito 2, 50 cursos ( 12,8 \% ) foram avaliados com o conceito 4, 19 cursos ( 4,9 \% ) foram avaliados com o conceito 1 e, sete cursos ( 1,8 \% ), com o conceito 5.
Nessa categoria, quatro cursos ( 1,0 \% ) ficaram sem conceito ( SC ).
Na rede pública, o conceito modal foi o 4, atribuído a 32 cursos, correspondentes a 52,5 \% dos 61 cursos da categoria.
Dos demais cursos participantes, 18 cursos ( 29,5 \% ) obtiveram o conceito 5, oito cursos ( 13,1 \% ) receberam o conceito 2 e três cursos ( 4,9 \% ), o conceito 3.
O conceito 1 foi atribuído a nenhum curso ( 0,0 \% ).
Nessa categoria administrativa, nenhum curso ( 0,0 \% ) ficou sem conceito ( SC ).
Dos 452 cursos participantes da modalidade presencial, o conceito 3 foi o valor modal, atribuído a 184 ( 40,7 \% ) cursos.
Dos demais cursos dessa modalidade de oferta, 138 cursos ( 30,5 \% ) receberam o conceito 2 ; 82 cursos ( 18,1 \% ) receberam o conceito 4 ; 25 cursos ( 5,5 \% ), o conceito 5 ; e 19 cursos ( 4,2 \% ), o conceito 1.
Nessa modalidade de oferta, quatro cursos ( 0,9 \% ) ficaram sem conceito ( SC ).
Na análise por região, observa -se, que, na região Norte, as instituições públicas participaram com quatro cursos ( 9,3 \% total regional ), dos quais nenhum ficou sem conceito ( SC ) ou recebeu o conceito 1.
Os conceitos modais, atribuídos a um curso cada, correspondem aos conceitos 2, 3, 4 e 5.
As instituições privadas participaram com 39 cursos na região Norte ( 90,7 \% do total regional ), e nenhum curso ficou sem conceito ( SC ).
O conceito modal nas instituições privadas na região foi o 2, atribuído a 20 cursos, o que corresponde a 51,3 \% dessa categoria na região.
Os conceitos 3, 1 e 5 foram atribuídos a 14, cinco e nenhum curso, respectivamente.
O conceito 4 foi atribuído a nenhum curso ( 0,0 \% ).
Dos cursos oferecidos na região Norte, 43 eram presenciais e nenhum ficou sem conceito ( SC ).
Dos 43 cursos oferecidos, 21 obtiveram o conceito 2 ( modal ) ; 15, o conceito 3 ; cinco, o conceito 1 ; um, o conceito 5 ; e um, o conceito 4.
Na região Nordeste, a rede privada concentrou 103 dos 122 cursos participantes, o equivalente a 84,4 \% do total da região.
Entre os cursos oferecidos por IES privadas no Nordeste, 43 cursos foram avaliados com conceito 3 ou o conceito 2, conceitos modais.
Nove cursos receberam o conceito 4 ; seis cursos, o conceito 1 ; e um curso obteve o conceito 5.
Nessa categoria administrativa e grande região, um curso ficou sem conceito ( SC ).
As instituições públicas da região Nordeste participaram com 19 cursos ( 15,6 \% do total da região ).
Desses, 13 ( 68,4 \% ) obtiveram conceito 4, o valor modal, seguidos por seis cursos, que receberam o conceito 5.
Os conceitos 1, 2 e 3 foram atribuídos a nenhum curso.
Nenhum curso ficou sem conceito ( SC ).
Dos cursos oferecidos na região Nordeste, todos ( 122 ) eram presenciais : um ficou sem conceito ( SC ), seis obtiveram o conceito 1 ; 43, o conceito 2 ; 43, o conceito 3 ; 22, o conceito 4 ; e sete, o conceito 5.
Na região Sudeste, a proporção de cursos da rede privada, 86,7 \%, correspondeu a 143 dos 165 cursos participantes.
Nessa combinação de categoria administrativa e grande região, o conceito modal foi o 3, atribuído a 72 cursos ; e um curso ficou sem conceito ( SC ).
Os demais foram avaliados com os seguintes conceitos : 2 ( 40 cursos ), 4 ( 23 cursos ), 5 ( cinco cursos ) e 1 ( dois cursos ).
Dos 22 cursos oferecidos por instituições públicas na região Sudeste, cinco obtiveram o conceito 5 ; 11, o conceito 4 ; um, o conceito 3 ; cinco obtiveram conceito 2 e nenhum obteve o conceito 1.
Nenhum curso ficou sem conceito ( SC ).
Ainda na região Sudeste, todos os cursos ( 100,0 \% ) foram oferecidos na modalidade presencial.
Para esses, a moda foi o conceito 3, com 73 cursos.
Os demais cursos presenciais da região foram \textbf{alocados} nos conceitos 2 ( 45 cursos ), 4 ( 34 cursos ), 5 ( dez cursos ) e 1 ( dois cursos ).
Um curso presenciais ficou sem conceito ( SC ).
As instituições privadas concentraram 66 dos 77 cursos participantes da região Sul, 85,7 \% do total regional.
Desses, 38 cursos obtiveram conceito 3, o conceito modal.
Nessa combinação de categoria administrativa e grande região, 13 receberam o conceito 4 ; 12 receberam o conceito 2 ; um, o conceito 5 ; e nenhum obteve o conceito 1.
Dois cursos ficaram sem conceito ( SC ).
As instituições públicas na região Sul participaram com 11 cursos ( 14,3 \% dos cursos da região ), aos quais foram atribuídos os conceitos 1 ( nenhum curso ), 2 ( nenhum curso ), 3 ( um curso ), 4 ( seis cursos ) e 5 ( quatro cursos ).
Nenhum curso da região Sul oferecido por IES públicas ficou sem conceito ( SC ).
Na região Sul, 77 dos cursos foram oferecidos na modalidade presencial.
Para esses, a moda foi o conceito 3, com 39 cursos.
Os demais cursos presenciais da região foram \textbf{alocados} nos conceitos 4 ( 19 cursos ), 2 ( 12 cursos ), 5 ( cinco cursos ) e 1 ( nenhum curso ).
Ficaram sem conceito ( SC ) dois cursos.
Na região Centro-Oeste, 40 dos 45 cursos participantes eram de IES privadas ( 88,9 \% em termos regionais ).
Desses, 15 cursos receberam conceito 2, o valor modal.
Dos demais cursos, 14 receberam o conceito 3, seis receberam o conceito 1, cinco receberam o conceito 4 e nenhum recebeu o conceito 5.
Nenhum ficou sem conceito ( SC ).
Quanto aos cinco cursos oferecidos por instituições públicas na região Centro-Oeste ( 11,1 \% do total regional ), nenhum ficou com conceito 3 ; um curso ficou sem conceito ( SC ) nessa região.
Também na região Centro-Oeste, todos os cursos foram ofertados na modalidade presencial.
Para esses, a moda foi o conceito 2, com 17 cursos ( 37,8 \% ).
Os demais cursos presenciais da região foram \textbf{alocados} nos conceitos 3 ( 14 cursos, 31,1 \% ), 4 ( seis cursos, 13,3 \% ), 1 ( seis cursos, 13,3 \% ) e 5 ( dois cursos, 4,4 \% ).
Nenhum curso presencial da região ficou sem conceito ( SC ).
A título de síntese, no Gráfico 4.2, apresenta -se a distribuição do Conceito Enade, segundo a categoria administrativa da IES.
Nesse gráfico, como dito anteriormente, poligonais mais à esquerda representam Conceito Enade mais baixos, e poligonais mais à direita, Conceito Enade mais altos.
É importante observar que, para cada poligonal presente nesse gráfico, a porcentagem constante no eixo y refere -se ao universo de cursos oferecidos por instituições públicas ( poligonal azul ) ou privadas ( poligonal verde ), respectivamente.
Os cursos em IES públicas ( poligonal azul ) apresentam uma poligonal mais à direita do que os em IES privadas ( poligonal verde ) e, consequentemente, uma distribuição de Conceito Enade com valores maiores. 4.3 CONCEITOS POR ORGANIZAÇÃO ACADÊMICA E GRANDE REGIÃO Na Tabela 4.3, encontra -se a distribuição dos conceitos atribuídos aos cursos participantes do Enade 2023, na área de Odontologia, por organização acadêmica, segundo as grandes regiões brasileiras. oferecido por CEFET/IF.
Essa distribuição corresponde a, respectivamente, 32,5 \%, 38,1 \%, 29,4 \% e 0,0 \% dos cursos.
De acordo com os dados apresentados na mesma tabela, de todos os 184 cursos avaliados com o conceito 3, 49 eram vinculados a universidades.
Os cursos vinculados a esse tipo de organização acadêmica tiveram o conceito 4 como conceito modal, atribuído a 50 cursos ( 34,0 \% ).
Os demais cursos ofertados por universidades receberam os conceitos 5 ( 23 cursos ), 2 ( 23 cursos ) e 1 ( um curso ).
Um curso ficou sem conceito ( SC ).
Dos cursos vinculados a centros universitários, o conceito modal foi o 3, atribuído a 78 cursos ( 45,3 \% ).
Um curso vinculado a esse tipo de organização acadêmica ficou sem conceito ( SC ), e os outros cursos receberam os conceitos 2 ( 63 cursos ), 4 ( 22 cursos ), 1 ( sete cursos ) e 5 ( um curso ).
Dos 133 cursos mantidos por faculdades, 57 ( 42,9 \% ) receberam conceito 3, o conceito modal.
O segundo conceito mais frequentemente atribuído aos cursos oferecidos por essa categoria de IES foi o conceito 2, recebido por 52 cursos.
Na sequência, 11 cursos obtiveram conceito 1 ; dez cursos, o conceito 4 ; e um curso, o conceito 5.
Dos cursos oferecidos em faculdades, dois cursos ficaram sem conceito ( SC ).
Considerando -se, separadamente, as regiões brasileiras, verifica -se que, na região Norte, as universidades participaram com oito cursos avaliados, aos quais foram atribuídos os seguintes conceitos : 1 a nenhum curso ; 2 a três cursos ; 3 a três cursos ; conceito 4 a um curso ; e 5 a um curso.
Nenhum curso ficou sem conceito ( SC ).
Os centros universitários da região Norte foram representados por 16 cursos, dos quais dois cursos obteve conceito 1 ; oito cursos, conceito 2, o valor modal ; seis cursos, conceito 3 ; nenhum curso, o conceito 4 ; e nenhum curso, o conceito 5.
Nenhum curso ficou sem conceito ( SC ).
As faculdades participaram com 19 cursos na região Norte, os quais receberam os conceitos 1 ( três cursos ), 2 ( dez cursos ), 3 ( seis cursos ), 4 ( nenhum curso ) e 5 ( nenhum curso ).
Nenhum curso ficou sem conceito ( SC ).
Na região Nordeste, as universidades participaram com 24 dos 122 cursos da área de Odontologia oferecidos.
O conceito modal foi 4, atribuído a 13 cursos.
Outros seis cursos obtiveram o conceito 5 ; quatro cursos, o conceito 3 ; um curso, o conceito 2 ; e nenhum curso, o conceito 1.
Nenhum curso ficou sem conceito ( SC ).
Os centros universitários contaram com 57 cursos participantes na região Nordeste, 24 dos quais receberam o conceito modal 3.
Para os demais cursos, foram atribuídos os conceitos 2 ( 22 cursos ), 4 ( sete cursos ), 1 ( três cursos ) e 5 ( um curso ).
Nenhum curso oferecido por centros universitários no Nordeste ficou sem conceito ( SC ).
As faculdades foram representadas por 41 cursos na região Nordeste, os quais receberam os conceitos 1 ( três cursos ), 2 ( 20 cursos ), 3 ( 15 cursos ), 4 ( dois cursos ) e 5 ( nenhum curso ).
Um curso ficou sem conceito ( SC ).
Na região Sudeste, as universidades concentraram 67 dos 165 cursos de Odontologia da região.
Dos cursos oferecidos em IES com este tipo de organização acadêmica, nessa região, o conceito modal foi 3, atribuído a 25 cursos.
Os demais cursos receberam os conceitos 4 ( 19 cursos ), 2 ( 13 cursos ), 5 ( nove cursos ) e 1 ( um curso ).
Nenhum curso ficou sem conceito ( SC ).
Os centros universitários participaram com 56 cursos na região Sudeste, dos quais 26 obtiveram o conceito modal 3.
O conceito 2 foi atribuído a 20 cursos ; o conceito 4 a dez cursos ; o conceito 5 a nenhum curso ; e o conceito 1 a nenhum curso.
Nenhum curso ficou sem conceito ( SC ) nessa combinação de organização acadêmica e região.
Já as faculdades foram representadas por 42 cursos na região Sudeste, dos quais um ficou sem conceito ( SC ).
O conceito modal 3 foi atribuído a 22 cursos.
Aos demais cursos foram atribuídos os conceitos 2 ( 12 cursos ), 4 ( cinco cursos ), 5 ( um curso ) e 1 ( um curso ).
Dos 77 cursos da região Sul, 38 eram mantidos por universidades, para os quais o conceito modal foi 4, atribuído a 15 cursos.
Os demais cursos receberam os conceitos 3 ( 13 cursos ), 5 ( cinco cursos ), 2 ( quatro cursos ) e 1 ( nenhum curso ).
Em IES com esse tipo de organização acadêmica, um curso da região Sul ficou sem conceito ( SC ).
Dos 26 cursos participantes de centros universitários da região Sul, um ficou sem conceito ( SC ), e 16 cursos receberam conceito 3, o valor modal.
Os conceitos restantes ficaram assim distribuídos : 2 ( sete cursos ), 4 ( dois cursos ), 5 ( nenhum curso ) e 1 ( nenhum curso ).
Quanto aos 13 cursos vinculados a faculdades na região Sul, dez cursos receberam o conceito 3, valor modal.
Os demais conceitos ficaram assim distribuídos : 4 ( dois cursos ), 2 ( um curso ), 5 ( nenhum curso ) e 1 ( nenhum curso ).
Dos 13 cursos, nenhum ficou sem conceito ( SC ).
Na região Centro-Oeste, 10 dos 45 cursos oferecidos eram mantidos por universidades.
Para esse tipo de organização acadêmica, o conceito modal foi 3, atribuído a quatro cursos.
Em relação às universidades, nenhum curso da região ficou sem conceito ( SC ).
Os demais cursos receberam os conceitos 5 ( dois cursos ), 4 ( dois cursos ), 2 ( dois cursos ) e 1 ( nenhum curso ).
Os centros universitários da região Centro-Oeste contaram com 17 cursos : seis cursos foram avaliados com o conceito 3 ; seis cursos, com o conceito 2 ; três cursos, com o conceito 4 ; dois cursos, com o conceito 1 ; e nenhum curso obteve o conceito 5.
Nenhum curso ficou sem conceito ( SC ).
Dos 18 cursos oferecidos em faculdades na região Centro-Oeste, nove cursos receberam o conceito 2 ( modal ).
Os conceitos restantes ficaram assim distribuídos : 3 ( quatro cursos ), 1 ( quatro cursos ), 4 ( um curso ) e 5 ( nenhum curso ).
Nessa região, nenhum curso oferecido em faculdades ficou sem conceito ( SC ) ou obteve o conceito 5.
De maneira semelhante à forma como se procedeu em relação à seção anterior, esta seção encerra -se com o Gráfico 4.3, que apresenta a distribuição do Conceito Enade dos cursos de Odontologia segundo a organização acadêmica da IES.
No Gráfico 4.3, apresenta -se a distribuição do Conceito Enade segundo a organização acadêmica da IES.
A leitura desse último gráfico é semelhante aos anteriores, seguindo -se a lógica de que poligonais deslocadas para a esquerda correspondem aos conceitos mais baixos ( 1, 2 e 3 ) e poligonais deslocadas para a direita correspondem a Conceitos Enade mais altos.
Nas universidades ( linha azul ), todos os cursos obtiveram conceito 5 ou abaixo disso, já que a poligonal cumulativa alcança 100 \% para o conceito 5.
Além disso, os cursos de universidades apresentam uma poligonal mais à direita do que as demais, denotando uma melhor distribuição de Conceitos Enade.
No outro extremo, encontram -se os cursos mantidos por faculdades ( linha vermelha ), apresentando as respectivas poligonais mais à esquerda e, consequentemente, uma distribuição de Conceitos Enade com os menores valores.
Os cursos mantidos por centros universitários ( linha verde ) apresenta poligonal mais central e, consequentemente, uma distribuição de Conceitos Enade com os valores intermediários.

% matched lemmas: alocar
\end{textsample}
